\section{Experimental Setup}
\label{sec:experiments}
We evaluate eleven configurations spanning five autonomous coding agents and nine frontier language models, designed to enable controlled ablation of the agent scaffold and the underlying model independently.

\paragraph{Agents.}
We select five coding agent frameworks: \textbf{OpenCode}~\cite{opencode_ai}, \textbf{Aider}~\cite{aider_chat}, \textbf{Codex}~\cite{chen2021codex}, \textbf{Claude~Code}~\cite{anthropic_claude_code}, and \textbf{Gemini}\footnote{\url{https://gemini.google.com/}}.
Each provides a distinct combination of scaffolding capabilities (tool integration, file-system access, context management). Codex, Claude~Code, and Gemini are proprietary systems tightly coupled to their respective model families. OpenCode and Aider are open-source frameworks that support multiple backend LLMs, enabling clean model ablations at fixed scaffold. All agents are evaluated under the same two-phase protocol described in Appendix~\ref{subsection: agent evaluation}: a read-only \emph{planning} phase in which the agent produces a structured optimization plan, followed by a writable \emph{execution} phase in which the agent implements only the changes specified in its own plan.

\paragraph{Language models.}
Codex is evaluated with its native \textbf{GPT-5}~\cite{openai_gpt5_page}. Claude~Code is evaluated with four Anthropic models: \textbf{Sonnet\,4.5}~\cite{anthropic_claude_sonnet_45}, \textbf{Sonnet\,4.6}, \textbf{Opus\,4.6}~\cite{anthropic_claude_opus_4_6_system_card_2026}, and \textbf{Opus\,4.7}. Gemini is evaluated with \textbf{2.5\,Pro} and \textbf{2.5\,Flash}. To enable a clean LLM ablation at fixed agent scaffold, we evaluate OpenCode with three models: \textbf{GPT-5.1\,Codex}, \textbf{GPT-5}, and \textbf{GPT-4.1}~\cite{openai_gpt41_2025}. Aider is evaluated with GPT-5, completing a GPT-5 triad (OpenCode, Aider, Codex) that isolates the agent scaffold at fixed model. All agents receive identical inputs: repository snapshot, deterministic host script, baseline CWV for both device profiles, and LCP/CLS/INP attribution signals, under reasoning and output token budget constraints detailed in \S\ref{sec:agent_protocol}.

\paragraph{Evaluation subset.}
From the full benchmark (\S\ref{sec:bench_sources}), we evaluate on a stratified subset of 300 websites spanning all 17 frameworks, selected to cover a range of baseline CWV profiles: sites with at least one Poor metric (requiring improvement), sites in the Needs Improvement range, and sites where all metrics are already Good (requiring restraint).
Each website is measured over 30 independent runs per device profile (mobile and desktop) as described in \S\ref{sec:cwv-evaluation}, yielding six metrics per configuration: LCP, INP, and CLS on each of mobile and desktop.
We aggregate per-site scores using the median across runs to obtain robust baselines and post-agent estimates.


\subsection{Metrics}

\subsubsection{Site Health}
\label{sec:site-health}
Comparing agents across six metrics on different scales LCP in milliseconds (range: 20-10{,}000+), INP in milliseconds (range: 0-2{,}000+), and CLS as a unitless ratio (range: 0-2+). We define a \emph{Site Health Score} that maps each metric to a common $[0,100]$ scale following the equation \ref{eq:health}:
\begin{equation}
\label{eq:health}
    H(v, m) = \begin{cases}
        100 & \text{if } v \leq G_m \\[3pt]
        100 - 50 \cdot \frac{v - G_m}{N_m - G_m} & \text{if } G_m < v \leq N_m \\[3pt]
        \max\!\Bigl(0,\; 50\cdot(1 - \frac{v - N_m}{N_m})\Bigr) & \text{if } v > N_m
    \end{cases}
\end{equation}
where $G_m$ and $N_m$ are the Good and Needs~Improvement thresholds for metric $m$ (e.g., $G_{\text{LCP}}=2500$\,ms, $N_{\text{LCP}}=4000$\,ms), the ranges and thresholds are consistent with the performance scoring tiers adopted by Adobe \footnote{\url{https://experienceleague.adobe.com/en/docs/experience-manager-learn/cloud-service/expert-resources/cloud-5/season-3/cloud5-lighthouse-score-optimization-part1}}. A score of 100 indicates the metric is at or below the Good threshold; 50 corresponds to the Needs~Improvement boundary; 0 indicates severe degradation. We report the \emph{Health Delta} the change in the mean of $H(v,m)$ across all six metrics from baseline to post-agent.